\subsection{Initial-visit MC-O-PI converges with uniform updates}
In the sampled algorithm, a single update can worsen the policy even though the mean field improves it. We show instead that each sufficiently late interval with a fixed target has a uniformly positive chance of ending in an improvement or of continuing forever. Boundedness lets us prove this statement before leaving a fixed ball. Finitely many target values then rule out infinitely many unsuccessful transitions.

In addition to the standing assumptions, impose full inertia \eqref{eq:greedy-pol-tie-breaking}, and sample the initial action uniformly within its state:
\begin{equation}\label{eq: uniform starts}
\sigma_k(s,a)=\sigma_k(s)\ge0,\qquad
\sum_s|\mathcal A(s)|\sigma_k(s)=1.
\end{equation}
Here $\sigma_k(s)$ is the probability of each pair at $s$; the state marginal is $|\mathcal A(s)|\sigma_k(s)$. Thus \eqref{eq: cumulative learning} reduces to $\sum_k\alpha_k\sigma_k(s)=\infty$ at every state. We also require the following weaker form of eventual nonincrease.
\begin{assumption}[Bounded future increases]\label{asp: additional conditions on learning rate}
There are deterministic $C_\alpha\ge1$ and $k_\alpha<\infty$ such that
\begin{equation}\label{eq:comparison}
\alpha_j\le C_\alpha\alpha_k\qquad(j\ge k\ge k_\alpha).
\end{equation}
\end{assumption}
Eventual nonincrease is the case $C_\alpha=1$. The condition also permits infinitely many increases, for example $\alpha_k=(2+(-1)^k)/(k+3)$ with $C_\alpha=3$. It imposes no lower bound on future steps relative to the current one.

\begin{theorem}\label{thm: convergence mcopi uniform}
Under the standing assumptions, initial-visit updates, uniform initial actions \eqref{eq: uniform starts}, full inertia \eqref{eq:greedy-pol-tie-breaking} and bounded future step increases \eqref{eq:comparison}, almost surely $(q_{\pi_k})$ changes only finitely many times and
\[
q_k\longrightarrow q_{\pi_*}.
\]
Every sufficiently late selected policy is optimal. If every state has a unique optimal action, $\pi_k$ is eventually constant.
\end{theorem}
\subsubsection{Relation to the combined stability--ODE method}
\label{sec: limitations of existing methods}
Our proof follows the combined stability--ODE strategy of \citet[Section~5.4]{Kushner2003StochasticApplications}: establish stability, then use repeated opportunities to follow improving mean-field behaviour. A difficulty here is that entries into a policy region $Q(\pi)$ may have arbitrarily small margins from its boundary. Recurrence to the region alone therefore does not supply the fixed positive distance used in a usual lock-in estimate; see, for example, the setup in \citet[Section~3]{Borkar2002OnApproximation}.

We address this difficulty by constructing favorable updates that create positive gaps against policy-worsening actions, then bounding the probability that subsequent noise closes those gaps. States complete this construction separately, and each gap is protected on its own accumulated-learning clock. Condition~\eqref{eq:comparison} compares earlier favorable steps with the completion-time margin and bounds future noise relative to that margin. It is sufficient for this argument; we do not claim it is necessary for convergence. An exit that improves the policy counts as progress, so the proof does not require permanent confinement to a single policy region.

The proof occupies the remainder of this section. The following probability estimate controls an estimated gap between two actions. Its dependence on a state's own accumulated learning, rather than elapsed iteration count, allows arbitrarily long waits between updates of that state.
\subsubsection{A barrier on a state's learning clock}
The same scalar weight appears in a gap's drift and variance. Keeping this weight allows arbitrarily long pauses. When it is zero, the noise vanishes and the drift cannot decrease the gap. No divergence of the accumulated weight is needed for the following local estimate.

\begin{lemma}[State-clock barrier]\label{lem:barrier}
Let $D_k$ be real and $\mathcal F'_k$-measurable, with integrable increments. Let $r\le\tau$ be stopping times, and let $\sigma_k(s)\in[0,1]$ be an $\mathcal F'_k$-measurable scalar weight for a fixed state $s$. In the application it will be the probability of each initial pair at $s$. The deterministic steps are in $(0,1]$, square summable, and satisfy $\alpha_j\le C_\alpha\alpha_k$ for all $j\ge k\ge k_\alpha$, with deterministic $C_\alpha\ge1$ and $k_\alpha$. Assume $r\ge k_\alpha$. Write
\[
m_k=\mathbb E[D_{k+1}-D_k\mid\mathcal F'_k],\qquad
Z_{k+1}=D_{k+1}-D_k-m_k.
\]
Suppose deterministic $\delta,V>0$ satisfy, on $\{r\le k<\tau,D_k\ge0\}$,
\begin{equation}\label{eq:barrierhyp}
m_k\ge\alpha_k\sigma_k(s)(4\delta-D_k),\qquad
\mathbb E[Z_{k+1}^2\mid\mathcal F'_k]\le V\alpha_k^2\sigma_k(s).
\end{equation}
Then, on $\{r<\infty,D_r>0\}$,
\begin{equation}\label{eq:barrier}
\mathbb P\bigl(\exists\text{ finite }n\in[r,\tau]:D_n\le0\mid\mathcal F'_r\bigr)
\le C_0\frac{C_\alpha\alpha_r}{\min\{D_r,\delta\}}+\frac{C_1}{\delta^2}\sum_{k=r}^{\infty}\alpha_k^2,
\end{equation}
where $C_0=\max\{1,20V(1+(3\delta)^{-1})\}$ and $C_1=16V$. The crossing event includes the update
ending at $\tau$.
\end{lemma}
\begin{proof}
\emph{Before reaching $\delta$.} Put $a_r=C_\alpha\alpha_r$, so every future step is at most $a_r$. Suppose $0<D_r<\delta$. If $D_r<a_r$, the first term in \eqref{eq:barrier} is already
at least one. Otherwise define
\[
\nu=\tau\wedge\inf\{n\ge r:D_n\le0\text{ or }D_n\ge\delta\},\qquad
M_n=\sum_{k=r}^{n-1}\mathbf1_{\{k<\nu\}}Z_{k+1},
\]
and the state clock
\begin{equation}\label{eq:clock}
A_s(r,n)=\sum_{k=r}^{n-1}\alpha_k\sigma_k(s).
\end{equation}
A first crossing at $n\le\tau$ before reaching $\delta$ forces
\begin{equation}\label{eq:crossing}
M_n\le-D_r-3\delta A_s(r,n).
\end{equation}
For $b>0$ known at $r$, set $T(b)=\inf\{n\ge r:A_s(r,n)\ge b\}$, possibly infinite. The bound $\alpha_k\le a_r$ controls the clock overshoot and future variance, at every finite horizon $n\ge r$, by
\[
A_s(r,T(b)\wedge n)\le b+a_r,\qquad
\sum_{k=r}^{T(b)\wedge n-1}\alpha_k^2\sigma_k(s)\le a_r(b+a_r).
\]
To justify the conditional estimate, first fix a deterministic restart $r$ and stop at $T(b)\wedge N$ for a finite horizon $N$. Stopping preserves zero conditional means because whether an increment is retained is known before that increment is drawn. Orthogonality and the variance bound give
\[
\mathbb E[M_{T(b)\wedge N}^2\mid\mathcal F'_r]
\le V a_r(b+a_r).
\]
The threshold-crossing proof given in Section~\ref{sec: replacement foundations} works conditionally too: multiply its expectations by the indicator of any event known at $r$. The tower property then gives the same inequality conditional on $\mathcal F'_r$. Quantities $b,z$ known at $r$ remain fixed in this calculation. In particular, the following bound holds (even the constant $1$ in place of $4$ suffices):
\begin{equation}\label{eq:doob}
\mathbb P\!\left(\sup_{r\le n\le T(b)}|M_n|\ge z\ \middle|\ \mathcal F'_r\right)
\le\frac{4Va_r(b+a_r)}{z^2},\qquad z>0.
\end{equation}
The supremum is over finite indices. First stop at a finite horizon, then increase the horizon;
an unattained supremum is handled by increasing thresholds to $z$ from below. For possibly unbounded $\mathcal F'_r$-measurable $b,z$, perform the calculation on $\{b\le m,\ 1/m\le z\le m\}$ and increase $m$; these events are known at $r$ and exhaust $\{b,z>0\}$. Random $r$ is treated on $\{r=j\}$
and then combined over integers $j$. Thus no finite bound on $T(b)$ or the waiting time has
been assumed.

Partition possible crossing times into $A_s(r,n)<D_r/(3\delta)$ and, for $j\ge1$,
\[
2^{j-1}D_r/(3\delta)\le A_s(r,n)<2^jD_r/(3\delta).
\]
Apply \eqref{eq:doob} with $(b,z)=(D_r/(3\delta),D_r)$ for the first range and $(b,z)=(2^jD_r/(3\delta),2^{j-1}D_r)$
for the others. Because $D_r\ge a_r$, the respective bounds are at most
\[
4V(1+(3\delta)^{-1})\frac{a_r}{D_r},\qquad
16V(1+(3\delta)^{-1})\frac{a_r}{D_r}2^{-j}.
\]
Summing and using $a_r=C_\alpha\alpha_r$ gives the first term in \eqref{eq:barrier}.

\emph{After reaching $\delta$.} Let $\rho=\inf\{n\ge r:D_n\ge\delta\}$, using $\rho=r$ if $D_r\ge\delta$. On $\{\rho<\infty,\rho\le\tau\}$,
restart at $\rho$ and stop at $\eta=\tau\wedge\inf\{n\ge\rho:D_n\le0\}$:
\[
\widetilde M_n=\sum_{k=\rho}^{n-1}\mathbf1_{\{k<\eta\}}Z_{k+1}.
\]
For $D_k\ge\delta$, \eqref{eq:barrierhyp} and $0\le\alpha_k\sigma_k(s)\le1$ imply $D_k+m_k\ge\delta$. If a first subsequent crossing is
at $n$, let $\ell<n$ be the last index with $D_\ell\ge\delta$. All subsequent gaps before $n$ are in $(0,\delta)$ and
have nonnegative mean increments, so
\[
D_n\ge\delta+\widetilde M_n-\widetilde M_\ell.
\]
A crossing forces $\sup_{n\ge\rho}|\widetilde M_n|\ge\delta/2$. Doob's inequality at $\rho$ bounds this probability by
$16V\delta^{-2}\sum_{k\ge\rho}\alpha_k^2$. Multiply by the $\mathcal F'_\rho$-measurable event $\{\rho<\infty,\rho\le\tau\}$, average back to $\mathcal F'_r$,
and use $\rho\ge r$. This proves the second term. The last-visit index $\ell$ is used only in a pathwise
inequality; we never condition at it. Both parts retain the increment $k=\tau-1$ and therefore
protect the gap through the endpoint itself.
\end{proof}

\subsubsection{Reference gaps and separately completed seeds}
For a reference policy $\pi$, define
\begin{equation}\label{eq:gaps}
\begin{gathered}
d_\pi(s,a)=q_\pi(s,\pi(s))-q_\pi(s,a),\qquad
\mathcal B_\pi=\{(s,a):d_\pi(s,a)>0\},\\
d_k(s,a)=q_k(s,\pi(s))-q_k(s,a).
\end{gathered}
\end{equation}
An action in $\mathcal B_\pi$ would worsen the reference action according to $q_\pi$. If these estimated gaps
are all positive, no greedy policy can select a bad action, so policy improvement applies. If
every $\mathcal B_\pi$ is empty, improvement applies between every pair of policies in both directions.
Their state values are equal; choosing a policy taking any prescribed action shows that every
action attains that same value. All policies are optimal with the same target, and Lemma~\ref{thm: bounded limit sets}
completes the theorem. Otherwise set
\begin{equation}\label{eq:constants}
d_* = \min_{\pi,\ (s,a)\in\mathcal B_\pi}d_\pi(s,a)>0,\qquad
\delta=d_*/4,\qquad h=g=\delta/2.
\end{equation}

Fix deterministic $K>B$ and restart $r$, and let $\zeta_r=\inf\{k\ge r:\|q_k\|_\infty>K\}$. For a
stopping time $t\ge r$, work on $\{t<\infty,t<\zeta_r\}$, fix $\pi=\pi_t$, and define
\begin{equation}\label{eq:stops}
\tau=\inf\{k>t:q_{\pi_k}\ne q_\pi\},\qquad \chi=\tau\wedge\zeta_r.
\end{equation}
These are stopping times. Return moments give finite second moments for every fixed-time
table by induction in \eqref{eq: update q with chi}. For $t\le k<\chi$, the action-value target is exactly $q_\pi$, even if the
selected policy label is different. Equal initial action probabilities give
\begin{align}
\mathbb E[d_{k+1}(s,a)-d_k(s,a)\mid\mathcal F'_k]
&=\alpha_k\sigma_k(s)(d_\pi(s,a)-d_k(s,a)),\label{eq:gapdrift}\\
\operatorname{Var}(d_{k+1}(s,a)-d_k(s,a)\mid\mathcal F'_k)
&\le V\alpha_k^2\sigma_k(s),\qquad V=2\{c_v+(B+K)^2\}.\label{eq:gapvariance}
\end{align}
Only the two coordinates forming the gap can change, their selection events are disjoint,
and each contributes at most $\alpha_k^2\sigma_k(s)\{c_v+(B+K)^2\}$ to the raw second moment. For each bad pair, $d_\pi(s,a)\ge4\delta$, so the barrier's drift and variance hypotheses hold through $\chi$. Its probability bound will be used once a positive starting gap has been obtained. All such moment estimates
are used before exit, not conditioned on future containment in the ball.

\paragraph{A favorable local update.}
Suppose $\pi_k(s)=\pi(s)$ and a current
bad gap is below $\delta$. Request one of the following action/return events:
\begin{enumerate}
\item If $q_\pi(s,\pi(s))-q_k(s,\pi(s))\ge d_*/2$, select $\pi(s)$ and require $G_{k+1}\ge q_\pi(s,\pi(s))-h$. The
reference coordinate increases by at least $3\delta\alpha_k/2$.
\item Otherwise select $a$ and require $G_{k+1}\le q_\pi(s,a)+h$. The identity
\[
q_k(s,a)-q_\pi(s,a)=q_k(s,\pi(s))-q_\pi(s,\pi(s))+d_\pi(s,a)-d_k(s,a)>\delta
\]
implies that the bad coordinate decreases by more than $g\alpha_k=\delta\alpha_k/2$.
\end{enumerate}
In either case no other gap against the reference action decreases, and the reference action
remains greedy. For centered $X$ with variance at most $c_v$, the one-sided Chebyshev inequality
gives
\[
\mathbb P(X\ge-h)\ge\frac{h^2}{c_v+h^2},\qquad
\mathbb P(X\le h)\ge\frac{h^2}{c_v+h^2}.
\]
For example, apply Markov's inequality to $(X+c_v/h)^2$ to bound $\mathbb P(X\ge h)$ and repeat for
$-X$; when $c_v=0$, $X=0$ almost surely. Conditional on selecting $s$, its initial action is
uniform, so every requested local update has conditional probability at least
\begin{equation}\label{eq:localprob}
p_0=\frac{h^2}{\max_s|\mathcal A(s)|\,(c_v+h^2)}>0.
\end{equation}
No probability of selecting the state enters this bound.

\begin{lemma}[Separately completed statewise seeds]\label{lem:seeds}
For full-inertia ties and any positive integer $N$, there is an adapted construction asking for at most $JN$
favorable local updates, starting at $t\ge k_\alpha$. Along successful requests, every unfinished state
retains its reference action, including at an endpoint reached by a favorable update. If state $s$
completes at time $T_s<\infty$, then
\begin{equation}\label{eq:seedmargin}
d_{T_s}(s,a)\ge\min\{\delta,gN\alpha_{T_s}/C_\alpha\}\qquad((s,a)\in\mathcal B_\pi).
\end{equation}
Each request, conditional on selecting its unfinished state, has probability at least $p_0$.
\end{lemma}
\begin{proof}
Order the bad actions at each state. Skip an action immediately if its gap is already
at least $\delta$; otherwise process it with at most $N$ favorable updates at that state, stopping
earlier if it reaches $\delta$. Process skips at $t$ and after every favorable update, before the next
sample. A state completes when its list is exhausted. Write $T_s$ for this time, using $T_s=t$ for
a state completed initially and $T_s=\infty$ if the construction stops before completion. These
are stopping times. Every selection of an unfinished state requests a local update; completed
states receive ordinary updates. Stop the construction at the first failed requirement or at $\chi$.

Under full inertia, raising the reference estimate or lowering a competitor retains the
reference action, since it remains greedy. At an unupdated state the table is unchanged, so
the action is also retained. Induction proves the assertion through every unfinished state's
wait.

The same retention conclusion also holds for fixed priorities: at time $t$ the reference action has highest priority among that state's
maximizers. Leaving the table unchanged, raising its estimate, or lowering a competitor
cannot create a new maximizer ahead of it in the priority order. Therefore that same action
continues to be selected at every unfinished state. Induct over all iterations, including those
updating other states and the final favorable update at completion. This proves the required
retention property directly; a fixed-priority rule is not being identified with full inertia on
arbitrary updates.

Every processed gap either reaches $\delta$ or receives $N$ increases of at least $g\alpha_k$. Later seed
updates at its state cannot lower it. Bounded future increases give $\alpha_k\ge\alpha_{T_s}/C_\alpha$
for every such update and proves \eqref{eq:seedmargin}. At most $J$ bad actions are processed, so there are
at most $JN$ requests, irrespective of waiting times. Their conditional probability bound is
\eqref{eq:localprob}.
\end{proof}

Completed states can be updated while other states wait. We therefore cannot condition
indiscriminately on all future successful behavior: the completed-state return laws must stay
unchanged. The following bounded change of measure conditions only the seed updates.

An auxiliary probability law is a way of reweighting possible runs for the proof. If a nonnegative random weight $L_\infty$ has expectation one, define
\[
\widehat{\mathbb P}(E)=\mathbb E[L_\infty\mathbf1_E]
\]
for every event $E$. Runs with larger weight become more likely under this law, and an event of original probability zero still has probability zero. This last property is denoted by $\widehat{\mathbb P}\ll\mathbb P$. We choose the weights so that requested updates always succeed, while the laws at completed states remain unchanged.

\begin{lemma}[Conditioning only the seed updates]\label{lem:measure}
There is an auxiliary law $\widehat{\mathbb P}\ll\mathbb P$
under which all active requirements of Lemma~\ref{lem:seeds} succeed. Conditional state selection is
unchanged. At completed states, and after the construction stops, the conditional action and
return law given the current history and selected state is unchanged. For every future event
$E$, on $\{t<\infty,t<\zeta_r\}$,
\begin{equation}\label{eq:transfer}
\mathbb P(E\mid\mathcal F'_t)\ge p_0^{JN}\widehat{\mathbb P}(E\mid\mathcal F'_t).
\end{equation}
\end{lemma}
\begin{proof}
Let $S_{k+1}$ be the state component of $I_{k+1}$. At an active unfinished state $s$, let
$p_k(s)$ be the conditional probability of its requirement given $\mathcal F'_k,S_{k+1}=s$. Then $p_k(s)\ge p_0$.
If the state has zero sampling probability, set $p_k(s)=p_0$; this value is never used on a
positive-probability draw.

Start $L_t=1$. At a requested update of $s$, multiply $L$ by the indicator of success divided
by $p_k(s)$; at other iterations multiply it by one. Every factor has conditional mean one given
$\mathcal F'_k,S_{k+1}$. Hence $(L_k)$ is a martingale, meaning $\mathbb E[L_{k+1}\mid\mathcal F'_k]=L_k$. At most $JN$ factors differ from one, even though their
times need not be bounded, so $0\le L_k\le p_0^{-JN}$. Since only finitely many factors differ from one, $L_k$ is eventually constant on each path. Bounded convergence says that expectations pass through an almost-sure limit when the variables share a deterministic finite bound. Indeed, if $X_n\to X$ almost surely and $|X_n|,|X|\le M$, then for every $\varepsilon>0$,
\[
\mathbb E|X_n-X|\le\varepsilon+2M\mathbb P(|X_n-X|>\varepsilon)\longrightarrow\varepsilon;
\]
letting $\varepsilon\downarrow0$ proves the claim. Apply this fact to $\mathbf1_A L_n$ for each event $A\in\mathcal F'_t$. Since the martingale identity gives $\mathbb E[\mathbf1_A L_n]=\mathbb P(A)$ for $n\ge t$, we obtain
\[
\mathbb E[\mathbf1_A L_\infty]=\mathbb P(A)\qquad(A\in\mathcal F'_t).
\]
The defining conditional-averaging property therefore yields $\mathbb E[L_\infty\mid\mathcal F'_t]=1$.
Define $\widehat{\mathbb P}(E)=\mathbb E[L_\infty\mathbf1_E]$ as above. For random $t$, use likelihood one before $t$, on $\{t=\infty\}$, and off $\{t<\zeta_r\}$,
and combine the constructions on $\{t=j\}$ over integers $j$. Each such event is known at $j$, so the fixed-start conditional calculation applies there. The law agrees with $\mathbb P$ on $\mathcal F'_t$
where the construction is active.

A failure makes $L_\infty=0$, so active requirements succeed almost surely under the auxiliary law. Here is why later requests do not bias a completed-state update. For any bounded random variable $Y$ determined by the next history $\mathcal F'_{k+1}$, the definition of weighted probabilities and the tower property give, on $\{L_k>0\}$,
\[
\widehat{\mathbb E}[Y\mid\mathcal F'_k]
=\frac{\mathbb E[L_\infty Y\mid\mathcal F'_k]}{L_k}
=\frac{\mathbb E[L_{k+1}Y\mid\mathcal F'_k]}{L_k}.
\]
For the first equality, put $\mathcal G=\mathcal F'_k$ and $H=\mathbb E[L_\infty Y\mid\mathcal G]/L_k$, setting $H=0$ where $L_k=0$. Since $\mathbb E[L_\infty\mid\mathcal G]=L_k$, the nonnegative weight $L_\infty$ is zero almost surely on $\{L_k=0\}$. For every $A\in\mathcal G$, the tower property therefore gives
\[
\widehat{\mathbb E}[\mathbf1_AH]
=\mathbb E[\mathbf1_A L_kH]
=\mathbb E[\mathbf1_A L_\infty Y]
=\widehat{\mathbb E}[\mathbf1_AY].
\]
This is precisely the conditional-averaging identity introduced in Section~\ref{sec: mcopi}, now under the weighted law. The second equality in the preceding conditional formula uses $\mathbb E[L_\infty\mid\mathcal F'_{k+1}]=L_{k+1}$, so future factors have averaged out. Write $R_{k+1}=L_{k+1}/L_k$ for the current factor. Since $\mathbb E[R_{k+1}\mid\mathcal F'_k,S_{k+1}]=1$, applying the formula to state indicators preserves the state-selection probabilities. Conditioning further on $S_{k+1}=s$ gives
\[
\widehat{\mathbb E}[Y\mid\mathcal F'_k,S_{k+1}=s]
=\mathbb E[R_{k+1}Y\mid\mathcal F'_k,S_{k+1}=s].
\]
At a completed state this factor is one, preserving its conditional action and return law. The same holds after the construction stops. Histories where $L_k=0$ have auxiliary probability zero, so no division by zero is needed. Absolute continuity preserves the almost-sure
algorithmic identities, including the full-inertia tie rule. Finally,
\[
\widehat{\mathbb P}(E\mid\mathcal F'_t)=\mathbb E[L_\infty\mathbf1_E\mid\mathcal F'_t]
\le p_0^{-JN}\mathbb P(E\mid\mathcal F'_t),
\]
which proves the bound. This auxiliary law is used only in the proof; the algorithm still
samples its original returns.
\end{proof}
